% ==============================================================================
% Plantilla Institucional de Trabajo Terminal — UPIIZ IPN
% Archivo: capitulos/protocolo/02-justificacion.tex
% ==============================================================================
% ESTRUCTURA NORMATIVA DE PROTOCOLO: CAPÍTULO 2 - JUSTIFICACIÓN
% ==============================================================================
% LÍMITE INSTITUCIONAL: Máximo 1 página.
%
% Directrices normativas (Anexo 1):
% - Explicar con claridad por qué se desarrollará el proyecto.
% - Mencionar los beneficios económicos, sociales y tecnológicos que aportará la solución.
% - Mencionar las ventajas respecto de dispositivos o sistemas semejantes, si existen.
% - En el párrafo final, incluir obligatoriamente al menos dos razones relacionadas con:
%   * desarrollo de conocimiento;
%   * aportaciones teóricas;
%   * solución de situaciones prácticas;
%   * aportaciones metodológicas.
% ==============================================================================

\chapter{Justificación}
\label{cap:proto-justificacion}

% [INSTRUCCIÓN PARA EL ALUMNO]:
% Redacte este capítulo en un máximo de una página respetando estrictamente los siguientes elementos:

[Explique de forma concreta por qué se desarrollará el proyecto, describiendo la necesidad detectada, el problema a atender y la pertinencia de formular esta propuesta de desarrollo].

\vspace{0.8em}
[Describa los beneficios económicos, sociales y tecnológicos derivados de la implementación del proyecto, tales como ahorro de costos, impacto en la calidad de vida, seguridad de los usuarios, incremento de eficiencia o modernización tecnológica].

\vspace{0.8em}
[Mencione las ventajas competitivas, técnicas u operativas que ofrece su propuesta respecto de dispositivos, máquinas, procesos o sistemas semejantes existentes en el mercado o en el ámbito académico, si existen].

\vspace{0.8em}
% PÁRRAFO FINAL OBLIGATORIO:
[En este párrafo final, sustente de forma explícita la pertinencia del proyecto incorporando al menos dos razones concretas relacionadas con: a) desarrollo de conocimiento, b) aportaciones teóricas, c) solución de situaciones prácticas, o d) aportaciones metodológicas].
